\subsection{Analysis of China's Soybean Import Trade}

Brazil, the United States, and Argentina are major soybean suppliers to China. In 2024, their soybean exports to China totaled approximately 103.5 million tonnes, valued at USD 45.95 billion. Brazil accounted for 70.1\% of export volume, followed by the United States (26.1\%) and Argentina (3.8\%), demonstrating China's high dependence on a few major suppliers.

From 2015 to 2024, the combined export volume of these three countries increased by 32.8\%, while their total export value rose by 53.9\%. Brazil's exports grew by 77.3\% in volume, with its market share increasing from 52.5\% to 70.1\%. In contrast, U.S. export volume declined slightly by 1.0\%, while Argentina's exports fell by 59.7\%, indicating a gradual shift toward Brazilian dominance.

Trade policy has also influenced import patterns. In 2018, amid the China-U.S. trade dispute, U.S. soybean exports to China dropped from 31.69 million tonnes to 8.24 million tonnes, while Brazil's market share surged to 85.5\%. Argentina's exports were additionally affected by domestic production shocks, particularly the severe drought in 2023.

Meanwhile, the weighted average unit value of soybean imports from these countries increased from approximately USD 383 per tonne in 2015 to USD 591 in 2022, before declining to USD 444 in 2024. The faster growth in export value than in volume highlights the importance of price fluctuations.

These trends suggest that China's soybean import market is shaped by trade costs, relative prices, supply capacity, and substitution among exporting countries. These factors should therefore be incorporated into the subsequent model to evaluate the effects of tariff policy changes on export volumes and market shares.